\documentclass{article}

%\begin{comment}
%offene Fragen: 
%- Pause hinter textbackslash
%- ûnd
%- Punkt unter Tilde
%- unterkringelte Sachen
%- mehrzeiliger Kommentar
%\end{comment}

\begin{document}



\textbf{Hello world}

% ich bin ein kommentar

% latex ignoriert whitespace, leerzeichen, absätze etc.
% ein neuer absatz entsteht durch eine leerzeile (mehrere absätze)
% mehrere leerzeilen: trotzdem nur ein absatz
% ein einfacher zeilenumbruch ohne neuen absatz macht man mit \\

Mit \% schreiben wir einen Kommentar. \\
Latex ignoriert Whitespace, Leerzeichen, Absätze etc.\\
Mehrere Leerzeilen: trotzdem nur ein Absatz.\\
Einen einfachen Zeilenumbruch ohne neuen Absatz macht man 
mit \textbackslash\textbackslash .\\

Escapen \\


Mit \textbackslash       escapen wir Zeichen wie \% . \\
Einen Backslash machen wir mit \textbackslash textbackslash. \\ 
Zeichen, die escaped werden müssen, sind: \% , \$ , \{\} , \textbackslash,
\& , \# , \_ , \^ und \~ . \\

Text in normalen "Anführungszeichen" sieht so aus. \\
Typographisch richtig werden die `Anführungszeichen'' so.
Mit jeweils zwei \`` Backticks und \' Apostrophen. \\

Bindestrich kurz - \\
Gedankenstrich lang mit zwei Bindestrichen -- \\
noch länger mit drei Bindestrichen --- \\

\begin{center}
    \large
    Dieser Text ist zentriert. \\
    \tiny
    mit begin und end center
\end{center} 


\noindent mit noindent: Absatz nicht eingerückt. \\
Verschiedene Schriftgrößen sind relativ, nicht absolut.
Es gibt normalsize und kleinere und größere Schriften. \\
Kleiner: tiny, scriptsize, footnotesize, small. \\
Größer: large, Large, LARGE, huge, HUGE. \\


\noindent Beispiele folgen: \\
\begin{tiny} This text is tiny. \end{tiny} \\
\begin{scriptsize} This text is scriptsize. \end{scriptsize} \\
\begin{footnotesize} This text is footnotesize. \end{footnotesize} \\
\begin{small} This text is small. \end{small} \\
\begin{normalsize} This text is normalsize. \end{normalsize} \\
\begin{large} This text is large. \end{large} \\
\begin{Large} This text is Large. \end{Large} \\
\begin{LARGE} This text is LARGE. \end{LARGE} \\
\begin{huge} This text is huge. \end{huge} \\
\begin{Huge} This text is Huge. \end{Huge} \\

newpage macht eine neue Seite

\newpage

\noindent 
\textit{textit steht für italic, 
der Text ist dann das Argument in geschweiften Klammern}\\
\textsl{textsl steht für slanted, also etwas schräg}\\
\textbf{textbf steht für boldface}\\
\texttt{texttt steht für typewriter, also wie Schreibmaschine}\\
\textsc{textsc steht für smallcaps, also Kapitälchen}\\
\emph{emph heißt: emphasised, also kursiv und betont im Gegensatz zu textit nur kursiv}\\


\end{document}